% ============================================================
% Question Bank: ch06-05 Cross-Sections of Three-Dimensional Figures
% Topic Title: Cross-Sections of Three-Dimensional Figures
% CCSS: 7.G.A.3
% Questions: 27 (15 MC + 6 SA + 6 Graphical)
% ============================================================

% --- Multiple Choice Questions (15) ---

\begin{practiceQuestion}{6-5-q01}{mc}
\begin{questionText}
A right rectangular prism is sliced with a plane parallel to its base. What shape is the resulting cross-section?
\end{questionText}
\choiceA{Triangle}
\choiceB{Circle}
\choiceC{Rectangle congruent to the base}
\choiceD{Trapezoid}
\correctAnswer{C}
\explanation{A plane parallel to the base of a rectangular prism creates a cross-section that is a rectangle identical in shape and size to the base.}
\end{practiceQuestion}

\begin{practiceQuestion}{6-5-q02}{mc}
\begin{questionText}
A square pyramid is sliced horizontally, parallel to its base, halfway between the base and the apex. What shape is the cross-section?
\end{questionText}
\choiceA{A triangle}
\choiceB{A square smaller than the base}
\choiceC{A square the same size as the base}
\choiceD{A pentagon}
\correctAnswer{B}
\explanation{A slice parallel to the square base of a pyramid creates a smaller square. As the cutting plane moves closer to the apex, the square cross-section becomes smaller.}
\end{practiceQuestion}

\begin{practiceQuestion}{6-5-q03}{mc}
\begin{questionText}
A cone is sliced with a vertical plane that passes through the apex. What shape is the cross-section?
\end{questionText}
\choiceA{Circle}
\choiceB{Triangle}
\choiceC{Rectangle}
\choiceD{Oval}
\correctAnswer{B}
\explanation{A vertical cut through the apex of a cone passes through the two slant edges and the diameter of the base, producing a triangular cross-section.}
\end{practiceQuestion}

\begin{practiceQuestion}{6-5-q04}{mc}
\begin{questionText}
A cylinder is sliced with a plane perpendicular to its base (a vertical cut through the center). What shape is the cross-section?
\end{questionText}
\choiceA{Circle}
\choiceB{Ellipse}
\choiceC{Rectangle}
\choiceD{Triangle}
\correctAnswer{C}
\explanation{A vertical cut through the center of a cylinder perpendicular to both circular bases produces a rectangle whose width equals the diameter and whose height equals the cylinder's height.}
\end{practiceQuestion}

\begin{practiceQuestion}{6-5-q05}{mc}
\begin{questionText}
Which 3D figure always produces a circular cross-section no matter how it is sliced?
\end{questionText}
\choiceA{Cube}
\choiceB{Rectangular prism}
\choiceC{Sphere}
\choiceD{Triangular prism}
\correctAnswer{C}
\explanation{Every possible planar slice through a sphere produces a circle. The size of the circle depends on where the plane cuts, but the shape is always circular.}
\end{practiceQuestion}

\begin{practiceQuestion}{6-5-q06}{mc}
\begin{questionText}
A right rectangular prism is sliced diagonally from one edge of the top face to the opposite edge of the bottom face. What is the shape of the cross-section?
\end{questionText}
\choiceA{Triangle}
\choiceB{Rectangle}
\choiceC{Pentagon}
\choiceD{Hexagon}
\correctAnswer{B}
\explanation{A diagonal slice through a rectangular prism from one edge of the top to the opposite edge of the bottom produces a rectangular cross-section. The rectangle is tilted relative to the base.}
\end{practiceQuestion}

\begin{practiceQuestion}{6-5-q07}{mc}
\begin{questionText}
A triangular pyramid (tetrahedron) is sliced with a plane parallel to its base. What shape is the cross-section?
\end{questionText}
\choiceA{A square}
\choiceB{A triangle smaller than the base}
\choiceC{A triangle congruent to the base}
\choiceD{A rectangle}
\correctAnswer{B}
\explanation{A plane parallel to the triangular base of a triangular pyramid creates a smaller triangle. The cross-section is similar to the base but reduced in size because the pyramid narrows toward the apex.}
\end{practiceQuestion}

\begin{practiceQuestion}{6-5-q08}{mc}
\begin{questionText}
A cube is sliced with a plane that passes through three edges that do not share a common vertex. What shape could the cross-section be?
\end{questionText}
\choiceA{Triangle}
\choiceB{Pentagon}
\choiceC{Hexagon}
\choiceD{Circle}
\correctAnswer{C}
\explanation{When a plane cuts through six faces of a cube (passing through one edge on each of three pairs of opposite faces), the cross-section is a hexagon. This is the classic hexagonal cross-section of a cube.}
\end{practiceQuestion}

\begin{practiceQuestion}{6-5-q09}{mc}
\begin{questionText}
Javier slices a right rectangular prism with a vertical cut perpendicular to the longest face. The prism is $8$ cm long, $5$ cm wide, and $3$ cm tall. If the cut is parallel to the $5$-by-$3$ face, what are the dimensions of the cross-section?
\end{questionText}
\choiceA{$8$ cm by $5$ cm}
\choiceB{$8$ cm by $3$ cm}
\choiceC{$5$ cm by $3$ cm}
\choiceD{$3$ cm by $3$ cm}
\correctAnswer{C}
\explanation{A vertical cut parallel to the $5$-by-$3$ face slices through the width and height of the prism, producing a $5$ cm by $3$ cm rectangle.}
\end{practiceQuestion}

\begin{practiceQuestion}{6-5-q10}{mc}
\begin{questionText}
A hexagonal prism is sliced with a plane parallel to its base. What is the shape of the cross-section?
\end{questionText}
\choiceA{Rectangle}
\choiceB{Triangle}
\choiceC{Hexagon}
\choiceD{Pentagon}
\correctAnswer{C}
\explanation{A cut parallel to the base of any prism produces a cross-section congruent to the base. Since the base is a hexagon, the cross-section is also a hexagon.}
\end{practiceQuestion}

\begin{practiceQuestion}{6-5-q11}{mc}
\begin{questionText}
Which cross-section is NOT possible when slicing a rectangular pyramid?
\end{questionText}
\choiceA{Triangle}
\choiceB{Rectangle}
\choiceC{Trapezoid}
\choiceD{Circle}
\correctAnswer{D}
\explanation{A rectangular pyramid has flat faces and straight edges. Slicing it can produce triangles (through the apex), rectangles (parallel to the base), or trapezoids (angled cuts). A circle is impossible because the figure has no curved surfaces.}
\end{practiceQuestion}

\begin{practiceQuestion}{6-5-q12}{mc}
\begin{questionText}
A right rectangular prism measuring $6$ cm $\times$ $4$ cm $\times$ $3$ cm is sliced horizontally at a height of $1$ cm above the base. What is the area of the cross-section?
\end{questionText}
\choiceA{$12$ cm$^2$}
\choiceB{$18$ cm$^2$}
\choiceC{$24$ cm$^2$}
\choiceD{$72$ cm$^2$}
\correctAnswer{C}
\explanation{A horizontal cut through a rectangular prism produces a rectangle with the same length and width as the base. The cross-section is $6 \times 4 = 24$ cm$^2$.}
\end{practiceQuestion}

\begin{practiceQuestion}{6-5-q13}{mc}
\begin{questionText}
A square pyramid has a base with side length $10$ cm. A horizontal slice is made parallel to the base at the midpoint of the pyramid's height. The cross-section is a square with side length $5$ cm. What is the area of the cross-section?
\end{questionText}
\choiceA{$10$ cm$^2$}
\choiceB{$25$ cm$^2$}
\choiceC{$50$ cm$^2$}
\choiceD{$100$ cm$^2$}
\correctAnswer{B}
\explanation{The cross-section at the midpoint is a square with side $5$ cm. Area $= 5 \times 5 = 25$ cm$^2$.}
\end{practiceQuestion}

\begin{practiceQuestion}{6-5-q14}{mc}
\begin{questionText}
A triangular prism has an equilateral triangle base with sides of $6$ cm and a height (depth) of $10$ cm. It is sliced vertically, parallel to one rectangular face. What shape is the cross-section?
\end{questionText}
\choiceA{Triangle}
\choiceB{Rectangle}
\choiceC{Trapezoid}
\choiceD{Hexagon}
\correctAnswer{B}
\explanation{A vertical cut through a triangular prism parallel to one of its rectangular faces produces a rectangle. The height of the rectangle is $10$ cm and its width depends on where the cut is made.}
\end{practiceQuestion}

\begin{practiceQuestion}{6-5-q15}{mc}
\begin{questionText}
Samira has a cone with a circular base. She makes a slanted cut that does not pass through the base and does not pass through the apex. What shape could the cross-section be?
\end{questionText}
\choiceA{Circle}
\choiceB{Ellipse}
\choiceC{Square}
\choiceD{Triangle}
\correctAnswer{B}
\explanation{A slanted cut through a cone that avoids both the apex and the base produces an ellipse. Only a cut parallel to the base gives a perfect circle; an angled cut that misses the apex creates an elongated oval (ellipse).}
\end{practiceQuestion}

% --- Short Answer Questions (6) ---

\begin{practiceQuestion}{6-5-q16}{sa}
\begin{questionText}
Name the shape of the cross-section when a plane slices a cylinder parallel to its base.
\end{questionText}
\correctAnswer{Circle}
\explanation{A plane parallel to the circular base of a cylinder cuts through the cylinder at a constant distance from the base, producing a circle congruent to the base.}
\end{practiceQuestion}

\begin{practiceQuestion}{6-5-q17}{sa}
\begin{questionText}
A right rectangular prism has dimensions $10$ cm $\times$ $7$ cm $\times$ $4$ cm. A horizontal cut is made through the prism. What is the perimeter of the cross-section?
\end{questionText}
\correctAnswer{$34$ cm}
\explanation{A horizontal slice of a rectangular prism produces a rectangle matching the base: $10$ cm by $7$ cm. Perimeter $= 2(10 + 7) = 34$ cm.}
\end{practiceQuestion}

\begin{practiceQuestion}{6-5-q18}{sa}
\begin{questionText}
A square pyramid with a base side length of $12$ cm is cut by a horizontal plane at one-third of its height from the base. The resulting cross-section is a square with side length $8$ cm. What is the area of the cross-section?
\end{questionText}
\correctAnswer{$64$ cm$^2$}
\explanation{The cross-section is a square with side $8$ cm. Area $= 8 \times 8 = 64$ cm$^2$.}
\end{practiceQuestion}

\begin{practiceQuestion}{6-5-q19}{sa}
\begin{questionText}
Describe the shape of the cross-section when a cone is sliced with a plane that passes vertically through its apex and the center of its base.
\end{questionText}
\correctAnswer{Triangle (specifically an isosceles triangle)}
\explanation{A vertical cut through the apex and center of the base of a cone slices along two slant heights and the diameter. The resulting cross-section is an isosceles triangle.}
\end{practiceQuestion}

\begin{practiceQuestion}{6-5-q20}{sa}
\begin{questionText}
A cube has edge length $6$ cm. It is sliced with a plane parallel to one face. What is the area of the cross-section?
\end{questionText}
\correctAnswer{$36$ cm$^2$}
\explanation{A cut parallel to a face of a cube produces a square cross-section with the same side length as the cube's edge: $6 \times 6 = 36$ cm$^2$.}
\end{practiceQuestion}

\begin{practiceQuestion}{6-5-q21}{sa}
\begin{questionText}
A rectangular prism is $9$ cm long, $5$ cm wide, and $4$ cm tall. A vertical cut is made parallel to the $9$-cm-by-$4$-cm face. What is the area of the resulting cross-section?
\end{questionText}
\correctAnswer{$36$ cm$^2$}
\explanation{A vertical cut parallel to the $9 \times 4$ face produces a $9$ cm by $4$ cm rectangle. Area $= 9 \times 4 = 36$ cm$^2$.}
\end{practiceQuestion}

% --- Graphical MC Questions (3) ---

\begin{practiceQuestion}{6-5-q22}{gmc}
\begin{questionText}
The diagram below shows a rectangular prism with a horizontal cutting plane. What shape is the cross-section?

\begin{center}
\begin{tikzpicture}[scale=0.8]
  % Back face
  \draw[thick,dashed] (1,1) -- (5,1) -- (5,3) -- (1,3) -- cycle;
  % Front face
  \draw[thick] (0,0) -- (4,0) -- (4,2) -- (0,2) -- cycle;
  % Connecting edges
  \draw[thick] (0,0) -- (1,1);
  \draw[thick] (4,0) -- (5,1);
  \draw[thick] (4,2) -- (5,3);
  \draw[thick] (0,2) -- (1,3);
  % Cutting plane (horizontal)
  \draw[thick,funRed,dashed] (0,1) -- (4,1) -- (5,2) -- (1,2) -- cycle;
  \node[funRed,right] at (5.2,1.5) {\small cut};
  % Labels
  \node[below] at (2,0) {$7$ cm};
  \node[left] at (0,1) {$5$ cm};
  \node[below right] at (4.7,0.3) {$4$ cm};
\end{tikzpicture}
\end{center}
\end{questionText}
\choiceA{Triangle}
\choiceB{Square}
\choiceC{Rectangle}
\choiceD{Trapezoid}
\correctAnswer{C}
\explanation{A horizontal plane through a rectangular prism creates a cross-section congruent to the base. The cross-section is a rectangle measuring $7$ cm by $4$ cm.}
\end{practiceQuestion}

\begin{practiceQuestion}{6-5-q23}{gmc}
\begin{questionText}
The diagram below shows a pyramid with a square base being sliced parallel to the base. What shape is the cross-section?

\begin{center}
\begin{tikzpicture}[scale=0.8]
  % Base
  \draw[thick] (0,0) -- (4,0) -- (5,1) -- (1,1) -- cycle;
  % Apex
  \coordinate (apex) at (2.5,4);
  % Edges to apex
  \draw[thick] (0,0) -- (apex);
  \draw[thick] (4,0) -- (apex);
  \draw[thick,dashed] (5,1) -- (apex);
  \draw[thick,dashed] (1,1) -- (apex);
  % Cutting plane
  \draw[thick,funRed,dashed] (0.8,1.3) -- (3.5,1.3) -- (4.2,2) -- (1.5,2) -- cycle;
  \node[funRed,right] at (4.4,1.6) {\small cut};
\end{tikzpicture}
\end{center}
\end{questionText}
\choiceA{Triangle}
\choiceB{A smaller square}
\choiceC{Pentagon}
\choiceD{Circle}
\correctAnswer{B}
\explanation{A horizontal slice parallel to the square base of a pyramid creates a smaller square. The cross-section is similar to the base but reduced in size because the pyramid tapers toward the apex.}
\end{practiceQuestion}

\begin{practiceQuestion}{6-5-q24}{gmc}
\begin{questionText}
The diagram shows a cone being sliced horizontally, parallel to the base. What shape is the cross-section?

\begin{center}
\begin{tikzpicture}[scale=0.8]
  % Cone
  \draw[thick] (-2,0) arc[start angle=180,end angle=360,x radius=2,y radius=0.6];
  \draw[thick,dashed] (-2,0) arc[start angle=180,end angle=0,x radius=2,y radius=0.6];
  \draw[thick] (-2,0) -- (0,4);
  \draw[thick] (2,0) -- (0,4);
  % Cutting plane
  \draw[thick,funRed,dashed] (-1.2,1.6) arc[start angle=180,end angle=360,x radius=1.2,y radius=0.35];
  \draw[thick,funRed] (-1.2,1.6) arc[start angle=180,end angle=0,x radius=1.2,y radius=0.35];
  \node[funRed,right] at (1.4,1.6) {\small cut};
\end{tikzpicture}
\end{center}
\end{questionText}
\choiceA{Triangle}
\choiceB{Rectangle}
\choiceC{Circle}
\choiceD{Oval (ellipse)}
\correctAnswer{C}
\explanation{A horizontal plane parallel to the circular base of a cone creates a circular cross-section. It is smaller than the base because the cone tapers upward toward the apex.}
\end{practiceQuestion}

% --- Graphical SA Questions (3) ---

\begin{practiceQuestion}{6-5-q25}{gsa}
\begin{questionText}
The diagram shows a triangular prism sliced by a plane parallel to its triangular base. Describe the shape and find the area of the cross-section if the triangular base has a base of $6$ cm and a height of $4$ cm.

\begin{center}
\begin{tikzpicture}[scale=0.7]
  % Front triangular face
  \draw[thick] (0,0) -- (4,0) -- (2,3) -- cycle;
  % Back triangular face
  \draw[thick,dashed] (1.5,1) -- (5.5,1) -- (3.5,4) -- cycle;
  % Connecting edges
  \draw[thick] (0,0) -- (1.5,1);
  \draw[thick] (4,0) -- (5.5,1);
  \draw[thick] (2,3) -- (3.5,4);
  % Cutting plane
  \draw[thick,funRed,dashed] (0.5,0.33) -- (3.5,0.33) -- (5,1.33) -- (2,1.33) -- cycle;
  \node[funRed,right] at (5.2,0.8) {\small cut};
  % Labels
  \node[below] at (2,0) {$6$ cm};
  \node[left] at (1.8,1.5) {$4$ cm};
\end{tikzpicture}
\end{center}
\end{questionText}
\correctAnswer{Triangle; $12$ cm$^2$}
\explanation{A plane parallel to the base of a prism produces a cross-section congruent to the base. The cross-section is a triangle with base $6$ cm and height $4$ cm. Area $= \frac{1}{2} \times 6 \times 4 = 12$ cm$^2$.}
\end{practiceQuestion}

\begin{practiceQuestion}{6-5-q26}{gsa}
\begin{questionText}
A rectangular prism with dimensions $8$ cm $\times$ $6$ cm $\times$ $5$ cm is sliced as shown. The cut passes vertically through the prism parallel to the $8$-cm-by-$5$-cm face. What is the area of the cross-section?

\begin{center}
\begin{tikzpicture}[scale=0.6]
  % Back face
  \draw[thick,dashed] (1.5,1.5) -- (9.5,1.5) -- (9.5,6.5) -- (1.5,6.5) -- cycle;
  % Front face
  \draw[thick] (0,0) -- (8,0) -- (8,5) -- (0,5) -- cycle;
  % Connecting edges
  \draw[thick] (0,0) -- (1.5,1.5);
  \draw[thick] (8,0) -- (9.5,1.5);
  \draw[thick] (8,5) -- (9.5,6.5);
  \draw[thick] (0,5) -- (1.5,6.5);
  % Cutting plane
  \draw[thick,funRed,fill=funRed!10,opacity=0.5] (0,0) -- (8,0) -- (8,5) -- (0,5) -- cycle;
  \node[funRed,below] at (4,-0.3) {\small cut};
  % Labels
  \node[below] at (4,0) {$8$ cm};
  \node[left] at (0,2.5) {$5$ cm};
  \node[below right] at (8.8,0.5) {$6$ cm};
\end{tikzpicture}
\end{center}
\end{questionText}
\correctAnswer{$40$ cm$^2$}
\explanation{A vertical cut parallel to the $8 \times 5$ face produces a rectangle of $8$ cm by $5$ cm. Area $= 8 \times 5 = 40$ cm$^2$.}
\end{practiceQuestion}

\begin{practiceQuestion}{6-5-q27}{gsa}
\begin{questionText}
The figure below shows a sphere being sliced by a plane. Describe the shape of the cross-section. If the sphere has a radius of $10$ cm and the plane passes through the center of the sphere, what is the area of the cross-section? (Use $\pi \approx 3.14$.)

\begin{center}
\begin{tikzpicture}[scale=0.6]
  % Sphere outline
  \draw[thick] (0,0) circle (3);
  % Equator (dashed back, solid front)
  \draw[thick,dashed] (-3,0) arc[start angle=180,end angle=360,x radius=3,y radius=0.8];
  \draw[thick] (-3,0) arc[start angle=180,end angle=0,x radius=3,y radius=0.8];
  % Cutting plane through center
  \draw[thick,funRed] (-3,0) arc[start angle=180,end angle=0,x radius=3,y radius=0.8];
  \draw[thick,funRed,dashed] (-3,0) arc[start angle=180,end angle=360,x radius=3,y radius=0.8];
  % Radius
  \draw[thick,<->] (0,0) -- (3,0);
  \node[above] at (1.5,0) {$10$ cm};
\end{tikzpicture}
\end{center}
\end{questionText}
\correctAnswer{Circle; $314$ cm$^2$}
\explanation{Any cross-section of a sphere is a circle. When the plane passes through the center, the cross-section is a great circle with the same radius as the sphere. Area $= \pi r^2 = 3.14 \times 10^2 = 314$ cm$^2$.}
\end{practiceQuestion}
